\chapter{SpanCraft 原型复盘：从概念堆积到单一回路}

SpanCraft 是一个桥梁建造教学原型。它的价值不在于已经证明了学习效果，而在于暴露了教学软件常见的设计错误。

\section*{第一版的问题}

第一版试图同时呈现材料、截面、成本、受力、规范和桥型。信息都“有用”，但玩家打开界面先面对参数和说明，尚未形成一个想解决的问题。结果更像把作业搬进网页。

\section*{重排原则}

修订时采用两个原则：先给魅力，再给负荷；一次只突出一个可感知的关系。玩家先搭、先试车、先看见变形或破坏，再逐步接触公式和约束。可视化放大只改变显示，不应伪装成真实变形；所有演示性处理都需要标明。

\section*{尚未完成的验证}

当前原型能够说明一种设计思路，但尚不能回答“是否提高了学习成绩”“效果能否迁移”“不同基础学生是否同样受益”。后续验证至少需要：明确学习目标、前后测或对照任务、过程日志、访谈，以及对错误概念的检查。

\section*{可复用的产品检查表}

\begin{itemize}
  \item 玩家每次操作会收到什么学科反馈？
  \item 不懂目标知识，是否也能靠反复点击过关？
  \item 显示放大、简化模型和真实计算是否清楚区分？
  \item 失败是否告诉玩家下一步可以检查哪里？
  \item 作品能否导出，供学生解释与教师评价？
\end{itemize}
